% !TeX root = closed-systems-from-comparison-completeness.tex
% ============================================================
\chapter*{Reader Guide}
\phantomsection
\addcontentsline{toc}{chapter}{Reader Guide}
\label{ch:reader-guide}
% ============================================================
This monograph is organized as a dependency chain.  A reader need not
read every technical construction on a first pass, but the order of
dependence matters: later smooth, field-theoretic, scalar, and
state-side statements are downstream from the comparison and transport
stack.

\section*{Primary path}
The shortest path through the main theorem architecture is
\begin{center}
\hyperref[ch:introduction]{Introduction}; \cref{ch:foundation,ch:bottom,ch:closure,ch:locus,ch:bridge,ch:interface,ch:einstein,ch:connection,ch:em,ch:matter,ch:numerical-closure,ch:dynamics-program}.
\end{center}
This path gives the admissibility boundary, quotient semantics, the
two-locus theorem, the transport obstruction bridge, the
intrinsic--smooth interface, the macroscopic compatibility law, the
connection-first reconstruction, the electromagnetic and scalar-sector
consequences, and the state-side reconstruction extension.

\section*{Foundational path}
Readers concerned mainly with the primitive status of the theory should
read
\begin{center}
\hyperref[ch:introduction]{Introduction}; \cref{ch:foundation,ch:bottom,ch:rotation,ch:closure}.
\end{center}
These chapters contain the finite-to-global admissibility boundary, the
derivation of rectangular completeness from intrinsicality and local
distinguishability, the quotient semantics, and the obstruction to
subsystem attribution.

\section*{Transport and geometry path}
Readers concerned mainly with holonomy, curvature, and smooth
realization should read
\begin{center}
\cref{ch:locus,ch:lift,ch:triangles,ch:descent_obstruction,ch:bridge,ch:christoffel,ch:interface,ch:connection}.
\end{center}
This path follows the transition from representative enrichment to
transport obstruction, then from quadratic transport defect to realized
curvature and connection-first geometry.

\section*{Physics-facing path}
Readers concerned mainly with the downstream physical sectors should
read
\begin{center}
\cref{ch:causality,ch:hilbert,ch:einstein,ch:connection,ch:em,ch:matter,ch:matter-numbers,ch:numerical-closure,ch:matter-realization,ch:dynamics-program}
and \hyperref[app:fusion]{Appendix B}.
\end{center}
This path should be read with the status distinctions in
\hyperref[ch:logical-status]{Logical Status and Dependency Ledger} in
view.  Some statements in these chapters are
intrinsic consequences of the comparison stack; others are
realization-level consequences after a smooth, Hilbert, phase, or
observer-side reading has been introduced.

\section*{Use of the front matter}
\hyperref[ch:notation-conventions]{Notation and Conventions} fixes the
recurring notation and conventions used throughout the book.
\hyperref[ch:standing-terminology]{Standing Terminology}
defines the recurring vocabulary at orientation level.  The formal
definitions remain in the chapters where they are used in proofs.
\hyperref[ch:logical-status]{Logical Status and Dependency Ledger}
records the logical status of the principal claims and the location at
which each enters the dependency chain.
